\documentclass[11pt,oneside]{article}

\usepackage[
  a4paper,
  top=2.4cm,
  bottom=2.1cm,
  left=2.3cm,
  right=2.3cm,
  headheight=15pt
]{geometry}

\usepackage[T1]{fontenc}
\usepackage[utf8]{inputenc}
\usepackage{lmodern}
\usepackage{microtype}
\usepackage{inconsolata}
\usepackage{amsmath,amssymb,mathtools,bm}
\usepackage{enumitem}
\usepackage{booktabs,array,multicol,longtable}
\usepackage{xcolor}
\usepackage[most]{tcolorbox}
\usepackage{titlesec}
\usepackage{fancyhdr}
\usepackage[hidelinks]{hyperref}

\setlength{\parindent}{0pt}
\setlength{\parskip}{0.55em}
\setlist[itemize]{topsep=3pt,itemsep=2pt,leftmargin=2.2em}
\setlist[enumerate]{topsep=3pt,itemsep=2pt,leftmargin=2.2em}

\definecolor{accent}{HTML}{1B3A5C}
\definecolor{q@frame}{HTML}{1A5276}
\definecolor{q@bg}{HTML}{F6FAFD}
\definecolor{key@frame}{HTML}{1E8449}
\definecolor{key@bg}{HTML}{F4FBF6}
\definecolor{note@frame}{HTML}{9A6A00}
\definecolor{note@bg}{HTML}{FEF9E7}

\tcbset{
  enhanced,
  breakable,
  arc=2.5mm,
  boxrule=0.5pt,
  left=4.5mm,
  right=4mm,
  top=2.5mm,
  bottom=2.5mm,
  fonttitle=\bfseries\small,
  sharp corners=west,
  parbox=false
}

\newtcolorbox{qbox}[1][]{
  colback=q@bg,
  colframe=q@frame,
  borderline west={3.5pt}{0pt}{q@frame},
  title={#1}
}

\newtcolorbox{keybox}[1][]{
  colback=key@bg,
  colframe=key@frame,
  borderline west={3.5pt}{0pt}{key@frame},
  title={#1}
}

\newtcolorbox{notebox}[1][]{
  colback=note@bg,
  colframe=note@frame,
  borderline west={3.5pt}{0pt}{note@frame},
  title={#1}
}

\titleformat{\section}
  {\normalfont\Large\bfseries\color{accent}}
  {}{0em}{}
  [{\vspace{1pt}\color{accent!45}\titlerule[0.45pt]}]
\titlespacing*{\section}{0pt}{3.2ex plus 1ex minus .2ex}{1.8ex plus .2ex}

\titleformat{\subsection}
  {\normalfont\normalsize\bfseries\color{accent!80!black}}
  {}{0em}{}
\titlespacing*{\subsection}{0pt}{2.2ex plus 1ex minus .2ex}{1.0ex plus .2ex}

\pagestyle{fancy}
\fancyhf{}
\fancyhead[L]{\small\color{gray!70}Old Exam Questions --- \coursetitle}
\fancyhead[R]{\small\color{gray!70}\thepage}
\renewcommand{\headrulewidth}{0.4pt}

\newcommand{\coursetitle}{Unsupervised Machine Learning}

\newcounter{question}
\newcommand{\questiontitle}[1]{\stepcounter{question}\begin{qbox}[Question \thequestion: #1]}
\newcommand{\closequestion}{\end{qbox}}
\newcommand{\repeatnote}{\textit{\color{gray!70}Asked multiple times.}}
\newcommand{\choice}{\item[$\square$] }
\newcommand{\answerline}[1]{\makebox[#1]{\hrulefill}}
\newcommand{\R}{\mathbb{R}}
\newcommand{\E}{\mathbb{E}}
\newcommand{\Prob}{\mathbb{P}}
\newcommand{\KL}[2]{D_{\mathrm{KL}}\!\left(#1\,\|\,#2\right)}
\newcommand{\dd}{\,\mathrm{d}}

\begin{document}

\thispagestyle{empty}
\vspace*{1.4cm}
\begin{center}
  {\color{accent}\rule{\linewidth}{1.8pt}}\par
  \vspace{0.65cm}
  {\fontsize{26}{32}\selectfont\bfseries\color{accent}Old Exam Questions\\[2pt] \coursetitle\par}
  \vspace{0.55cm}
  {\color{accent!35}\rule{\linewidth}{0.5pt}}\par
  \vspace{0.55cm}
  {\large Stefan Eder\par}
  \vspace{0.15cm}
  {\small\color{gray!70}\today\par}
\end{center}

\vspace{0.9cm}
\begin{notebox}[How this file is organized]
This file is separate from the course summary and is meant as a clean practice set.
Exact duplicates were removed. When essentially the same question appeared in multiple old exams,
that is marked with \repeatnote\ inside the question box.
The wording was lightly normalized where the scans were messy, but the mathematical content and answer choices were kept intact.
\end{notebox}

\tableofcontents
\newpage

\section{Probability, Information Theory, and Estimation}

\questiontitle{Cumulants for independent uniform variables}
Assume you have two independent random variables $X$ and $Y$ which are uniformly distributed on $[-1,1]$.
As in the lecture, let $\kappa_n(X)$ denote the $n$-th order cumulant of a random variable $X$.
Compute the following quantities and answer the last classification question.
Use the hint
\[
\kappa_4(X)=\E(X^4)-3\bigl(\E(X^2)\bigr)^2
\]
for centered $X$.

\begin{itemize}
  \item $\kappa_4(X)=\answerline{3.0cm}$
  \item $\kappa_4(2X)=\answerline{3.0cm}$
  \item $\kappa_4(X+Y)=\answerline{3.0cm}$
  \item $X$ is sub-Gaussian / super-Gaussian: \answerline{4.0cm}
\end{itemize}
\closequestion

\questiontitle{Cumulants for symmetric Bernoulli variables}
Assume you have two independent, discrete random variables $X$ and $Y$, both following independent symmetric Bernoulli distributions taking only the values $+1$ and $-1$, each with probability $1/2$.
As in the lecture, let $\kappa_n(X)$ denote the $n$-th order cumulant of a random variable $X$.
Compute the quantities below.
Use the hint
\[
\kappa_3(X)=\E(X^3), \qquad \kappa_4(X)=\E(X^4)-3\bigl(\E(X^2)\bigr)^2
\]
for centered $X$.

\begin{itemize}
  \item $\kappa_3(X)=\answerline{3.0cm}$
  \item $\kappa_4(X)=\answerline{3.0cm}$
  \item $\kappa_4(3X)/\kappa_4(X)=\answerline{3.0cm}$
  \item $\kappa_4(X+Y)/\kappa_4(X)=\answerline{3.0cm}$
\end{itemize}
\closequestion

\questiontitle{Eigenvectors, covariance matrix, and a scalar expression}
You are given a dataset with $n$ samples and three features.
The covariance matrix $C$ has eigenvalues $(\lambda_1,\lambda_2,\lambda_3)=(5,3,1)$
and associated normalized eigenvectors $(v_1,v_2,v_3)$ with $\lVert v_i\rVert=1$.
Compute
\[
A=(2v_1^\top+v_2^\top+3v_3^\top)\,C\,(2v_2+v_3).
\]
Enter the final scalar value:
\[
A=\answerline{3.0cm}
\]
\closequestion

\questiontitle{KL divergence for continuous uniform distributions}
\repeatnote

Assume $Y$ and $Z$ are uniformly distributed on $[0,\theta_Y]$ and $[0,\theta_Z]$, respectively:
\[
p_Y(x)=\begin{cases}
\frac{1}{\theta_Y}, & x\in[0,\theta_Y],\\
0, & \text{else,}
\end{cases}
\qquad
p_Z(x)=\begin{cases}
\frac{1}{\theta_Z}, & x\in[0,\theta_Z],\\
0, & \text{else.}
\end{cases}
\]
For $\theta_Y=e^5$ and $\theta_Z=e^{10}$, compute
\[
D_{\mathrm{KL}}(p_Y\|p_Z)=\int_{-\infty}^{\infty} p_Y(x)\log\!\left(\frac{p_Y(x)}{p_Z(x)}\right)\dd x.
\]
Use the natural logarithm and round to two decimal places.
\[
D_{\mathrm{KL}}(p_Y\|p_Z)=\answerline{3.0cm}
\]
\closequestion

\questiontitle{Outer product matrix}
Given two vectors
\[
x=\begin{pmatrix}1\\2\\3\end{pmatrix},
\qquad
y=\begin{pmatrix}4\\5\\6\end{pmatrix},
\]
compute the outer product matrix
\[
A=xy^\top.
\]
Enter the integer values of the third row of $A$:
\[
A_{31}=\answerline{2.0cm}\,\qquad A_{32}=\answerline{2.0cm}\,\qquad A_{33}=\answerline{2.0cm}.
\]
\closequestion

\questiontitle{Marginal probabilities, entropies, and mutual information --- variant 1}
\repeatnote

Assume you are given two binary random variables $X$ and $Y$ taking only the values $0$ and $1$ where their joint probability distribution is specified by
\[
\begin{array}{c|cc}
 & Y=0 & Y=1\\
\hline
X=0 & 0.25 & 0.125\\
X=1 & 0.25 & 0.375
\end{array}
\]
This means, for example, that $P(X=0,Y=0)=0.25$.
Compute the quantities below using the natural logarithm.
Round the final values to two decimal places.
\begin{itemize}
  \item $P(X=1)=\answerline{2.5cm}$
  \item $P(Y=1)=\answerline{2.5cm}$
  \item $H(X)=\answerline{2.5cm}$
  \item $H(Y)=\answerline{2.5cm}$
  \item $I(X;Y)=\answerline{2.5cm}$
\end{itemize}
\closequestion

\questiontitle{Marginal probabilities, entropies, and mutual information --- variant 2}
\repeatnote

Assume you are given two binary random variables $X$ and $Y$ taking only the values $0$ and $1$ where their joint probability distribution is specified by
\[
\begin{array}{c|cc}
 & Y=0 & Y=1\\
\hline
X=0 & \tfrac13 & \tfrac13\\
X=1 & 0 & \tfrac13
\end{array}
\]
This means, for example, that $P(X=0,Y=0)=\tfrac13$.
Compute the quantities below using the natural logarithm.
Round the final values to two decimal places.
\begin{itemize}
  \item $P(Y=0)=\answerline{2.5cm}$
  \item $H(X)=\answerline{2.5cm}$
  \item $H(Y)=\answerline{2.5cm}$
  \item $I(X;Y)=\answerline{2.5cm}$
\end{itemize}
\closequestion

\questiontitle{Marginal probabilities, entropies, and mutual information --- variant 3}
Assume you are given two binary random variables $X$ and $Y$ taking only the values $0$ and $1$ where their joint probability distribution is specified by
\[
\begin{array}{c|cc}
 & Y=0 & Y=1\\
\hline
X=0 & \tfrac{3}{10} & \tfrac{4}{10}\\
X=1 & \tfrac{2}{10} & \tfrac{1}{10}
\end{array}
\]
Compute the following quantities using the natural logarithm.
Round the final values to two decimal places.
\begin{itemize}
  \item $P(X=0)=\answerline{2.5cm}$
  \item $P(Y=0)=\answerline{2.5cm}$
  \item $H(X)=\answerline{2.5cm}$
  \item $H(Y)=\answerline{2.5cm}$
  \item $I(X;Y)=\answerline{2.5cm}$
\end{itemize}
\closequestion

\questiontitle{Parameter Estimation and Maximum Likelihood}
\repeatnote

Which statements about Parameter Estimation and Maximum Likelihood are true?
Multiple answers are possible.
\begin{itemize}
  \choice If the data are distributed i.i.d., the likelihood can be written as a product of individual probabilities over the data samples.
  \choice The likelihood and the log-likelihood are maximized for the same parameter value, and at that parameter value $\hat w$, their values coincide: $\mathcal{L}(\hat w)=\ln \mathcal{L}(\hat w)$.
  \choice Parameter estimation and the maximum likelihood approach rest on the assumption that the model class is not known and estimated from the data.
  \choice The maximum likelihood estimator is asymptotically unbiased and efficient, but not asymptotically consistent.
  \choice The mean squared error can be written as the sum of the variance of the estimator and the squared bias of the estimator.
  \choice An estimator $\hat w$ is unbiased if its expected value equals the true parameter: $\E(\hat w)=w$.
  \choice The likelihood $\mathcal{L}(\{x\};w)$ is the probability that the model with parameter $w$ produces the data $\{x\}$.
  \choice The mean squared error is defined as $\mathrm{mse}(\hat w)=(\hat w-w)^\top(\hat w-w)$.
\end{itemize}
\closequestion

\questiontitle{Dropdown selection task on estimation theory}
\repeatnote

Fill in the correct choices.

For any parameter $w$, an estimator $\hat w$ is \textit{consistent / efficient / unbiased} if $\E(\hat w)=w$.
The \textit{expected value / bias / mean squared error / variance} of $\hat w$ can be decomposed as a sum of its
\textit{expected value / bias / mean squared error / variance} and its squared
\textit{expected value / bias / mean squared error / variance}.

Let $X_1,\dots,X_n$ be random variables drawn independently from a distribution $p(\cdot;w)$.
Consider
\[
\frac{1}{n+1}\sum_{i=1}^n X_i
\]
as an estimator for the mean of the distribution $p(\cdot;w)$.
This estimator is
\textit{neither unbiased nor consistent / unbiased but not consistent / biased but consistent / both unbiased and consistent}.

Now consider $X_1$ as an estimator for the mean of the distribution $p(\cdot;w)$.
This estimator is
\textit{neither unbiased nor consistent / unbiased but not consistent / biased but consistent / both unbiased and consistent}.

Let $\{x_1,\dots,x_n\}$ be a sample drawn independently from the distribution $p(\cdot;w)$.
The function
\[
\prod_{i=1}^n p(x_i;w)
\]
is called the \textit{probability / loss / density / likelihood} function.
Since $p(\cdot;w)$ is usually a \textit{probability / loss / density / likelihood} function, the probability measure of the set $\{x_1,\dots,x_n\}$ is in fact zero.

The Cram\'er--Rao lower bound relates the \textit{expected value / bias / mean squared error / variance} of the estimator $\hat w$ to the inverse of its Fisher information.
An estimator that reaches the Cram\'er--Rao lower bound is said to be \textit{consistent / efficient / unbiased}.
\closequestion

\questiontitle{Unsupervised Machine Learning}
Which statements about Unsupervised Machine Learning are true?
Multiple answers are possible.
\begin{itemize}
  \choice Underfitting means that a model is adapted to special training characteristics such as noisy measurements or outliers.
  \choice Projection methods project data into a higher-dimensional space while keeping their neighborhoods.
  \choice The model selection procedure is called ``training'' or ``learning''.
  \choice Examples of unsupervised methods are: PCA, ICA, Factor Analysis, k-means clustering, and mixture models.
  \choice In parametric models, every parameter (vector) realization represents a model.
  \choice $k$-nearest neighbors is a nonparametric model.
  \choice Generative models do not aim at modeling the data creation process.
  \choice Unsupervised Machine Learning algorithms learn patterns from unlabeled data.
\end{itemize}
\closequestion

\questiontitle{Expectation Maximization (EM) and Maximum Entropy --- version A}
Which statements about Expectation Maximization (EM) and Maximum Entropy are correct?
Multiple answers are possible.
\begin{itemize}
  \choice The EM algorithm can be applied to hidden Markov models and mixture models.
  \choice In the M-step, the model parameters are updated.
  \choice The E-step increases the upper bound on the log-likelihood.
  \choice In EM, one deals with problems that have hidden variables as well as model parameters.
  \choice The entropy of a probability distribution is always larger equal $0$ and smaller equal $1$.
  \choice Hidden state models in general have analytical solutions for the likelihood.
  \choice In the E-step, the hidden variables parameters are updated.
  \choice The Maximum Entropy approach requires minimal prior assumptions to assign a-priori probabilities.
\end{itemize}
\closequestion

\questiontitle{Expectation Maximization (EM) and Maximum Entropy --- version B}
Which statements about Expectation Maximization (EM) and Maximum Entropy are correct?
Multiple answers are possible.
\begin{itemize}
  \choice In the M-step, the hidden variables are updated.
  \choice Maximum Entropy requires minimal prior assumptions to assign a-priori probabilities.
  \choice The E-step increases the upper bound on the log-likelihood.
  \choice In the E-step, the model parameters are updated.
  \choice The EM algorithm can be applied to hidden Markov models and mixture models.
  \choice The entropy of a probability distribution is always larger equal $0$ and smaller equal $1$.
  \choice The E-step increases the lower bound on the log-likelihood.
  \choice In EM, one deals with problems that have hidden variables as well as model parameters.
\end{itemize}
\closequestion

\section{PCA, Kernel PCA, and ICA}

\questiontitle{Principal Component Analysis --- statements, version A}
Which statements about Principal Component Analysis (PCA) are correct?
Multiple answers are possible.
\begin{itemize}
  \choice The first principal component corresponds to the largest eigenvalue in the eigendecomposition of the sample covariance matrix.
  \choice The singular values from the singular value decomposition of the data matrix $X$ are equal to the eigenvalues of the eigendecomposition of $X^\top X$.
  \choice The PCA solution is unique, i.e. there is only one solution.
  \choice Eigenvectors can be computed iteratively with Oja's rule.
  \choice The sample variance of the $k$-th projection is equal to the $k$-th eigenvalue of the sample covariance matrix.
  \choice The matrix of eigenvectors $U$ is an orthogonal matrix.
  \choice The sample covariance matrix is symmetric and positive semidefinite.
  \choice PCA transforms data into a new coordinate system such that the data has the largest mean value along the first coordinate, the second largest mean value along the second coordinate, and so on.
\end{itemize}
\closequestion

\questiontitle{Principal Component Analysis --- statements, version B}
We keep the notation of the lecture slides: $X$ is the data matrix and $U$ is the matrix of eigenvectors of $X^\top X$.
Which statements about PCA are true?
Multiple answers are possible.
\begin{itemize}
  \choice PCA transforms data into a new coordinate system such that the data has the largest variance along the first coordinate, the second largest variance along the second coordinate, and so on.
  \choice The first principal component corresponds to the largest eigenvalue in the eigendecomposition of the sample covariance matrix.
  \choice If there are $n$ data samples, each of which is $d$-dimensional, the matrix $U$ has dimension $n\times d$.
  \choice The sample covariance matrix is orthogonal and positive semidefinite.
  \choice The singular values from the singular value decomposition of the data matrix $X$ are equal to the square root of the eigenvalues of the eigendecomposition of $X^\top X$.
  \choice The eigenvalues obtained from the eigendecomposition of $X^\top X$ are all positive and sum up to $1$.
  \choice PCA is commonly used to visualize data by downprojecting it to a smaller number of dimensions, keeping the most relevant information.
  \choice The matrix $U$ is a symmetric matrix.
\end{itemize}
\closequestion

\questiontitle{Picture-based interpretation question: PCA, ICA, and EM}
The original exam contained three picture-based subquestions.
State which interpretation is more likely in each case.
\begin{enumerate}
  \item PCA and ICA are applied to a 2D cloud of points, leading to outcomes A and B.
  Which picture is more likely the result of PCA and which one the result of ICA?
  \vspace{0.4em}

  PCA: \answerline{3.0cm}\qquad ICA: \answerline{3.0cm}

  \item Two plots show a dataset in a coordinate system spanned by principal components.
  Which is more likely: picture 1 is shown in $\mathrm{PC2}$ vs. $\mathrm{PC1}$ and picture 2 in $\mathrm{PC3}$ vs. $\mathrm{PC1}$, or vice versa?
  \vspace{0.4em}

  Answer: \answerline{7.0cm}

  \item The left figure shows an original dataset, while figures $a$ and $b$ show results of running EM for a mixture of two Gaussians.
  Which is more likely: figure $a$ is after $5$ iterations and figure $b$ after $20$ iterations, or the other way around?
  \vspace{0.4em}

  Answer: \answerline{7.0cm}
\end{enumerate}
\closequestion

\questiontitle{PCA computation on a rank-1 matrix --- projected second sample}
\repeatnote

You are given four data samples, each of dimension $3$.
The data matrix reads
\[
X=\begin{pmatrix}
1&1&1\\
2&2&2\\
3&3&3\\
4&4&4
\end{pmatrix}.
\]
The eigenvectors of $X^\top X$ are
\[
u_1=\frac{1}{\sqrt3}(1,1,1)^\top,
\qquad
u_2=\frac{1}{\sqrt2}(-1,0,1)^\top,
\qquad
u_3=\frac{1}{\sqrt2}(-1,1,0)^\top.
\]
They are already normalized and already correctly ordered corresponding to eigenvalues in descending order.
Compute the PCA projection matrix $Y$ and enter the components of the projected second feature vector $y_2$.
Round to two decimal places.
\[
y_2^{(1)}=\answerline{2.5cm}\qquad y_2^{(2)}=\answerline{2.5cm}\qquad y_2^{(3)}=\answerline{2.5cm}
\]
\closequestion

\questiontitle{PCA computation on a centered matrix --- projected third sample}
\repeatnote

You are given four data samples, each of dimension $3$.
The data matrix reads
\[
X=\begin{pmatrix}
-1&0&1\\
-2&0&2\\
-3&0&3\\
-4&0&4
\end{pmatrix}.
\]
The eigenvectors of the covariance matrix are
\[
u_1=\frac{1}{\sqrt2}(-1,0,1)^\top,
\qquad
u_2=\frac{1}{\sqrt2}(1,0,1)^\top,
\qquad
u_3=(0,1,0)^\top.
\]
They are already normalized and already correctly ordered corresponding to eigenvalues in descending order.
Compute the PCA projection matrix $Y$ and enter the components of the projected third sample $y_3$.
Round to two decimal places.
\[
y_3^{(1)}=\answerline{2.5cm}\qquad y_3^{(2)}=\answerline{2.5cm}\qquad y_3^{(3)}=\answerline{2.5cm}
\]
\closequestion

\questiontitle{Kernel PCA}
Which statements about Kernel PCA are correct?
Use the notation from the lecture: data points $x_1,\dots,x_n\in\mathcal X=\R^m$, kernel $k:\mathcal X\times\mathcal X\to\R$, feature map $\phi:\mathcal X\to\mathcal H$, feature covariance matrix
\[
C=\frac1n\sum_{i=1}^n\phi(x_i)\phi(x_i)^\top,
\]
and Gram matrix $K$ with entries $K_{ij}=\phi(x_i)^\top\phi(x_j)$.
Multiple answers are possible.
\begin{itemize}
  \choice A disadvantage of Kernel PCA is that $\phi$ always needs to be computed explicitly, which is computationally costly.
  \choice In general, any nonnegative and symmetric function $k:\mathcal X\times\mathcal X\to\R$ defines a kernel.
  \choice To compute the projection of a data point $x$ onto the eigenvectors of $K$, you need to center $x$ and $K$, as well as to normalize the eigenvectors of $K$.
  \choice The main task of Kernel PCA is to solve the eigenvalue problem of $C$ in feature space. This can be achieved by solving the eigenvalue problem of $K$ in a Euclidean space.
  \choice In general, $\phi$ is a nonlinear map that maps into a higher-dimensional feature space $\mathcal H$.
  \choice Linear PCA can be seen as a special case of Kernel PCA by using $k(x,y)=x^\top y$ as the associated kernel function.
  \choice For Kernel PCA, it is essential to center $K$ before computing its eigendecomposition, since the eigenvalues of $K$ can in general be negative, which leads to inconclusive results.
  \choice $K$ and $C$ always have the same dimensionality and the same rank.
\end{itemize}
\closequestion

\questiontitle{Independent Component Analysis --- version A}
Which statements about Independent Component Analysis (ICA) are correct?
Multiple answers are possible.
\begin{itemize}
  \choice Maximizing the negentropy maximizes the mutual information.
  \choice In general, ICA components are neither ranked nor orthogonal.
  \choice Non-Gaussianity can be measured by the second cumulant.
  \choice Criteria for independence are mutual information between components and non-Gaussianity of components.
  \choice The ICA solution is unique.
  \choice For ICA, one needs at least as many sources as observations.
  \choice Infomax maximizes the mutual information between components.
  \choice FastICA relies on whitening and rotating the data.
\end{itemize}
\closequestion

\questiontitle{Independent Component Analysis --- version B}
Which statements about Independent Component Analysis (ICA) are correct?
Multiple answers are possible.
\begin{itemize}
  \choice Independence can be quantified by the mutual information between components or the non-Gaussianity of components.
  \choice Fast ICA uses whitening and rotating the data.
  \choice Maximizing the negentropy minimizes the mutual information.
  \choice For non-Gaussian sources, the ICA solution is unique up to permutation and scaling.
  \choice Non-Gaussianity can be measured by the second and third cumulant.
  \choice For ICA, the number of sources must be smaller than, or equal to, the number of observations.
  \choice ICA components are ranked and orthogonal.
  \choice Infomax maximizes the mutual information between components.
\end{itemize}
\closequestion

\questiontitle{Independent Component Analysis --- version C}
Which statements about Independent Component Analysis (ICA) are correct?
Multiple answers are possible.
\begin{itemize}
  \choice For ICA, one needs more sources than observations.
  \choice Maximizing the negentropy maximizes the mutual information.
  \choice Criteria for independence are mutual information between components and non-Gaussianity of components.
  \choice The ICA solution is unique.
  \choice Fast ICA relies on whitening and rotating the data.
  \choice Infomax maximizes the mutual information between components.
  \choice Non-Gaussianity can be measured by the second cumulant.
  \choice In general, ICA components are neither ranked nor orthogonal.
\end{itemize}
\closequestion

\section{Factor Analysis}

\questiontitle{Factor Analysis: low-rank covariance part}
In factor analysis, part of the covariance function is formed by
\[
C_U = UU^\top = \sum_{i=1}^{\ell} u_i u_i^\top,
\]
where $U\in\R^{m\times \ell}$, the vectors $u_i$ are the columns of $U$, and $\ell<m$.
Which constraints does this impose on $C_U$?
Multiple answers are possible.
\begin{itemize}
  \choice The maximum rank of the matrix is $\ell$.
  \choice The number of zero eigenvalues of $C_U$ is greater or equal to $m-\ell$.
  \choice Any covariance matrix in $m$ dimensions can be written in the form above.
  \choice $C_U$ is singular.
\end{itemize}
\closequestion

\questiontitle{Factor Analysis: dimensions of the objects}
We keep the notation from the lecture.
Here $n$ is the number of data samples, $m$ is the dimension of each sample, $\ell$ is the number of factors,
$y$ is the factor vector, $U$ is the factor loading matrix, $\Psi$ is the diagonal noise covariance matrix, and $\varepsilon$ is the noise vector.
Which statements are correct?
Multiple answers are possible.
\begin{itemize}
  \choice $\Psi$ has shape $m\times m$.
  \choice $\Psi$ has shape $\ell\times \ell$.
  \choice $\Psi$ has shape $n\times \ell$.
  \choice $U$ has shape $m\times \ell$.
  \choice $U$ has shape $m\times n$.
  \choice $U$ has shape $\ell\times m$.
  \choice $\varepsilon$ has dimension $m$.
  \choice $\varepsilon$ has dimension $\ell$.
  \choice $\varepsilon$ has dimension $n$.
  \choice $y$ has dimension $m$.
  \choice $y$ has dimension $\ell$.
  \choice $y$ has dimension $n$.
\end{itemize}
\closequestion

\questiontitle{Factor Analysis: compute the first row of the covariance matrix}
You are given the factor loading matrix
\[
U=\begin{pmatrix}
1&2\\
3&4\\
5&6
\end{pmatrix}
\]
and the diagonal noise covariance matrix
\[
\Psi=\begin{pmatrix}
5&0&0\\
0&5&0\\
0&0&5
\end{pmatrix}.
\]
Compute the first row of the covariance matrix
\[
C=UU^\top+\Psi.
\]
Enter integer values:
\[
C_{11}=\answerline{2.0cm}\qquad C_{12}=\answerline{2.0cm}\qquad C_{13}=\answerline{2.0cm}
\]
\closequestion

\questiontitle{Factor Analysis: conceptual statements}
Which statements about Factor Analysis (FA) are correct?
Multiple answers are possible.
\begin{itemize}
  \choice The factors are unique only up to orthogonal transformations (rotations).
  \choice Correlations between observations can only be explained by factors since the noise covariance matrix $\Psi$ is diagonal.
  \choice FA describes the variability of observations in terms of latent variables and noise.
  \choice The factor loading matrix $U$ describes the independent part of the covariance of the observations.
  \choice FA decomposes the covariance matrix of the data as $C=UU^\top+\Psi$.
  \choice FA assumes Gaussian noise and non-Gaussian latent variables.
  \choice Factors explain the variance of the observations, while the correlation is explained by noise.
\end{itemize}
\closequestion

\section{Clustering, Mixtures, and Projection Methods}

\questiontitle{$k$-means: one iteration, dataset A}
Assume you want to do $k$-means on the following 2-dimensional dataset:
\[
X=\{(1,0),(0,1),(-1,0),(0,-1),(2,0),(2,3),(5,0)\}.
\]
The initial cluster centers are $\mu=(0,1)$ and $\nu=(2,2)$.
Compute the coordinates of the new cluster centers $\mu^{\mathrm{new}}$ and $\nu^{\mathrm{new}}$ after one iteration of $k$-means.
\[
\mu^{\mathrm{new}}=(\answerline{1.7cm},\answerline{1.7cm}),
\qquad
\nu^{\mathrm{new}}=(\answerline{1.7cm},\answerline{1.7cm})
\]
\closequestion

\questiontitle{$k$-means and affinity propagation}
Which statements about $k$-means and affinity propagation are correct?
Multiple answers are possible.
\begin{itemize}
  \choice The $k$-means objective is to minimize the within-cluster sum of squared deviations.
  \choice In affinity propagation: the larger the preference, the smaller is the number of cluster centers.
  \choice In affinity propagation: the larger the damping factor, the faster the convergence.
  \choice In affinity propagation, cluster centers are data points.
  \choice In affinity propagation, the similarity between two points is typically chosen as their negative squared distance.
  \choice $k$-means always converges to the global optimum.
  \choice The parameter $k$ in $k$-means is automatically chosen by the algorithm to best fulfill the optimization objective.
  \choice $k$-means is sensitive to initialization of the cluster centers.
\end{itemize}
\closequestion

\questiontitle{$k$-means: one iteration, dataset B}
Assume you want to do $k$-means on the following 2-dimensional dataset:
\[
X=\{(1,0),(2,-1),(1,-2),(0,-1),(2,0),(2,1),(3,0),(3,1)\}.
\]
The initial cluster centers are $\mu=(0,-2)$ and $\nu=(4,1)$.
Compute the coordinates of the new cluster centers $\mu^{\mathrm{new}}$ and $\nu^{\mathrm{new}}$ after one iteration of $k$-means.
\[
\mu^{\mathrm{new}}=(\answerline{1.7cm},\answerline{1.7cm}),
\qquad
\nu^{\mathrm{new}}=(\answerline{1.7cm},\answerline{1.7cm})
\]
\closequestion

\questiontitle{$k$-means: one iteration, dataset C}
Assume you want to do $k$-means on the following 2-dimensional dataset:
\[
X=\{(1,-1),(0,0),(-1,-1),(0,-2),(2,0),(2,3),(5,0)\}.
\]
The initial cluster centers are $\mu=(0,0)$ and $\nu=(2,1)$.
Compute the coordinates of the new cluster centers $\mu^{\mathrm{new}}$ and $\nu^{\mathrm{new}}$ after one iteration of $k$-means.
\[
\mu^{\mathrm{new}}=(\answerline{1.7cm},\answerline{1.7cm}),
\qquad
\nu^{\mathrm{new}}=(\answerline{1.7cm},\answerline{1.7cm})
\]
\closequestion

\questiontitle{$t$-SNE}
Which statements about $t$-distributed Stochastic Neighborhood Embedding are correct?
Multiple answers are possible.
\begin{itemize}
  \choice It models the neighborhood in both observation and embedding space by a Gaussian distribution.
  \choice It minimizes the Kullback--Leibler divergence between distributions in observation space and embedding space.
  \choice It models the neighborhood of a data point in observation space proportional to a Student's $t$-distribution.
  \choice It is usually trained by gradient descent.
\end{itemize}
\closequestion

\questiontitle{Scaling and Projection Methods --- version A}
Which statements about Scaling and Projection Methods are correct?
Multiple answers are possible.
\begin{itemize}
  \choice Isomap uses a geodesic distance matrix, whose eigenvectors are the coordinates of the projected space.
  \choice The objective in $t$-SNE is minimization of the KL divergence between the similarities of the high-dimensional inputs and the low-dimensional mapped outputs.
  \choice In locally linear embedding (LLE), each data point is described as a linear combination of its neighbors.
  \choice The objective of scaling is to project data into a higher-dimensional space, e.g. to allow better model selection.
  \choice In self-organizing maps the objective is to minimize an energy (error) function.
  \choice Multidimensional scaling (MDS) keeps the distances between data points as well as possible.
  \choice LLE in general performs linear mappings to realize neighborhood-preserving embeddings.
  \choice In $t$-SNE, the similarity between data points $i$ and $j$ is given by their joint probability.
\end{itemize}
\closequestion

\questiontitle{Projection and scaling methods --- version B}
Which statements about projection and scaling methods are correct?
Multiple answers are possible.
\begin{itemize}
  \choice Non-negative Matrix Factorization (NMF) performs a multiplicative decomposition of a positive semidefinite matrix into two positive semidefinite matrices.
  \choice Locally Linear Embedding (LLE) is a linear embedding method that aims to preserve the neighborhood of each data point.
  \choice Multidimensional Scaling (MDS) aims to find a low-dimensional representation of the data that preserves pairwise Euclidean distances.
  \choice Stochastic Neighbor Embedding (SNE) defines the similarity of two data points as a conditional probability while $t$-SNE defines it as a joint probability.
  \choice Isomap measures distances between data points as the length of the shortest path along a neighborhood graph.
\end{itemize}
\closequestion

\questiontitle{Scaling and projection methods --- version C}
Tick the correct statements about scaling and projection methods.
We collect the original data points in a matrix $X$ and the mapped points in a matrix $Y$.
Multiple answers are possible.
\begin{itemize}
  \choice The $t$-distribution has heavier tails than the Gaussian distribution; thus in $t$-SNE the joint distributions $q_{ij}$ between $y_i$ and $y_j$ are computed using a $t$-distribution in order to alleviate the crowding problem.
  \choice Isomap tries to approximate the geodesic distance between two points on a manifold by constructing a neighborhood graph.
  \choice In SNE and $t$-SNE, the perplexity term is computed using the entropy of the similarities $q_{ij}$ of the mapped data points $y_i$ and $y_j$.
  \choice In SNE and $t$-SNE, the objective is the mean squared distance between $x_i$ and $y_i$.
  \choice Locally linear embedding and multidimensional scaling are probabilistic methods that rely on minimizing the KL divergence between the distributions of $X$ and $Y$.
  \choice In SNE and $t$-SNE, the computation of the similarity $p_{ij}$ between $x_i$ and $x_j$ relies on assuming a Gaussian distribution centered at $x_i$.
\end{itemize}
\closequestion

\questiontitle{Scaling and Projection Methods --- version D}
\repeatnote

Which statements about Scaling and Projection Methods are correct?
Multiple answers are possible.
\begin{itemize}
  \choice Multidimensional scaling (MDS) keeps the distances between data points as well as possible.
  \choice The objective of scaling is to project data into a higher-dimensional space, e.g. to allow better model selection.
  \choice Isomap uses a geodesic distance matrix, whose eigenvectors are the coordinates of the projected space.
  \choice In $t$-SNE, the similarity between data points $i$ and $j$ is given by their joint probability.
  \choice In self-organizing maps the objective is to minimize an energy (error) function.
  \choice The objective in $t$-SNE is minimization of the KL divergence between the similarities of the high-dimensional inputs and the low-dimensional mapped outputs.
  \choice In locally linear embedding (LLE), each data point is described as a linear combination of its neighbors.
  \choice LLE in general performs linear mappings to realize neighborhood-preserving embeddings.
\end{itemize}
\closequestion

\questiontitle{Non-negative Matrix Factorization (NMF)}
As in the lecture, let
\[
X=YU^\top,
\qquad Y\in\R^{n\times \ell},
\qquad U\in\R^{m\times \ell},
\]
with non-negativity constraints.
Which statements about NMF are correct?
Multiple answers are possible.
\begin{itemize}
  \choice NMF objectives are the KL divergence and the Euclidean distance.
  \choice NMF is always exactly solvable.
  \choice $X\in\R^{n\times m}$.
  \choice The non-negativity constraints make the representation of the observations purely additive.
  \choice The outer product representation reads $X=\sum_{k=1}^{\ell} y_k u_k^\top$.
  \choice NMF factorizes any positive or negative definite matrix into two positive definite matrices.
\end{itemize}
\closequestion

\questiontitle{Mixturetown bread question}
In Mixturetown there are two bakers, Albert $(A)$ and Bibi $(B)$.
Their loaf weights are modeled as Gaussian distributions
\[
f(x;\mu_A,\sigma_A),\qquad f(x;\mu_B,\sigma_B),
\]
with
\[
\mu_A=1\text{ kg},\quad \sigma_A=0.1\text{ kg},\qquad
\mu_B=1.5\text{ kg},\quad \sigma_B=0.2\text{ kg}.
\]
Seventy-five percent of all loaves are made by Bibi.
You observe a loaf of weight $1.1$ kg.
Determine the probability that Bibi made this loaf.
Use the given density values
\[
f(1.1;1,0.1)=2.41971,
\qquad
f(1.1;1.5,0.2)=0.26995.
\]
Round to two decimal places.
\[
P(B\mid x=1.1)=\answerline{2.5cm}
\]
\closequestion

\questiontitle{Mixture model posterior probability}
\repeatnote

Assume the following case:
\begin{itemize}
  \item The prior probability of component $j$ is $0.1$.
  \item The probability for observing a data sample $x$ given the total parameter set $\theta$ is $0.2$.
  \item The likelihood for observing a data sample $x$, given component $j$ and corresponding parameter vector, is $0.5$.
\end{itemize}
Compute the posterior probability of component $j$.
Round to two decimal places.
\[
P(j\mid x)=\answerline{2.5cm}
\]
\closequestion

\questiontitle{Mixture model: infer the likelihood from posterior, prior, and evidence}
Assume the following case:
\begin{itemize}
  \item The prior probability of component $j$ is $0.1$.
  \item The posterior probability of component $j$ is $0.3$.
  \item The probability for observing a data sample $x$ given the total parameter set $\theta$ is $0.2$.
\end{itemize}
Compute the likelihood for observing a data sample $x$ given component $j$ and the corresponding parameter vector.
Round to two decimal places.
\[
P(x\mid j)=\answerline{2.5cm}
\]
\closequestion

\questiontitle{Mixture of Gaussians: posterior of the first Gaussian}
The data generation process in a mixture of Gaussians model is as follows:
first choose one Gaussian according to its weight in the mixture, then draw a sample from that Gaussian.
You have fitted a mixture of two Gaussians to 1D data:
\[
p(x)=\alpha_1 f(x;\mu_1,\sigma_1)+\alpha_2 f(x;\mu_2,\sigma_2),
\]
with
\[
\alpha_1=0.3,\quad \alpha_2=0.7,\quad \mu_1=0,\quad \sigma_1=1,\quad \mu_2=3,\quad \sigma_2=1.
\]
Relevant densities are
\[
f(1;0,1)=0.242,
\qquad
f(1;3,1)=0.054.
\]
Calculate the probability that the datapoint $x=1$ has been produced by the first Gaussian.
Round to two decimal places.
\[
P(G_1\mid x=1)=\answerline{2.5cm}
\]
\closequestion

\questiontitle{KL divergence for continuous uniform distributions --- variant 2}
Assume $Y$ and $Z$ are uniformly distributed on $[0,\theta_Y]$ and $[0,\theta_Z]$, respectively.
For $\theta_Y=e^4$ and $\theta_Z=e^{10}$, compute
\[
D_{\mathrm{KL}}(p_Y\|p_Z)=\int_{-\infty}^{\infty} p_Y(x)\log\!\left(\frac{p_Y(x)}{p_Z(x)}\right)\dd x.
\]
Use the natural logarithm and round to two decimal places.
\[
D_{\mathrm{KL}}(p_Y\|p_Z)=\answerline{2.5cm}
\]
\closequestion

\questiontitle{KL divergence for two discrete distributions}
You are given two probability distributions $p$ and $q$:
\[
p(x=0)=0.8,\quad p(x=1)=0.2,
\qquad
q(x=0)=0.5,\quad q(x=1)=0.5.
\]
Compute the KL divergence
\[
D_{\mathrm{KL}}(p\|q).
\]
Use the natural logarithm and round to two decimal places.
\[
D_{\mathrm{KL}}(p\|q)=\answerline{2.5cm}
\]
\closequestion

\questiontitle{Clustering --- version A}
Which statements about Clustering are correct?
Multiple answers are possible.
\begin{itemize}
  \choice In affinity propagation, the cluster centers are always data points.
  \choice Clustering (or Cluster Analysis) is both a regression and a classification technique.
  \choice In both $k$-means and fuzzy $k$-means, the membership distribution is always given by a binary random variable.
  \choice In mixture models, an observation is assigned to the component with the largest prior distribution.
  \choice Top-down clustering relies on building the minimal spanning tree and then in each step removes the largest of the remaining edges.
  \choice In a mixture of Gaussians, the maximum likelihood and the maximum posterior always have the same value.
  \choice The agglomerative hierarchical clustering algorithm repeatedly fuses the two maximally distant clusters.
  \choice Similarity-based clustering does not require objects to be represented via feature vectors.
\end{itemize}
\closequestion

\questiontitle{Clustering --- version B}
Which statements about Clustering are correct?
Multiple answers are possible.
\begin{itemize}
  \choice $k$-means is faster than MoG and not sensitive to initialization.
  \choice Both $k$-means and MoG assign an observation to a single cluster only.
  \choice The mixing coefficients in a Mixture of Gaussians model are represented by the posteriors.
  \choice The maximum a posteriori estimation is a way to avoid near-zero variances and singular solutions in the maximum likelihood estimation of mixture models.
  \choice Given an observation $x$, a component $j$, and a parameter vector $\theta$, the posterior $p(j\mid x,\theta)$ describes how likely it is that $x$ was produced by $j$.
  \choice The EM algorithm for MoG tries to maximize the likelihood with respect to the cluster centers, the covariance matrices, and the mixing coefficients.
  \choice The $k$-means algorithm is a simplification of the Mixture of Gaussians model.
  \choice Clustering is an unsupervised technique that aims at finding separate regions of high data density called clusters.
\end{itemize}
\closequestion

\questiontitle{Hierarchical clustering: single linkage vs. complete linkage}
You are given $8$ datapoints in $\R^2$:
\[
\{(0,1),(1,1),(2,1),(8,1),(0,-1),(1,-1),(2,-2),(8,-1)\}.
\]
Agglomerative hierarchical clustering is run until only two clusters remain, once with single linkage and once with complete linkage.
Two resulting clusterings are shown in the exam figure.
Check the correct statements.
\begin{itemize}
  \choice Clustering 1 has been obtained by single linkage.
  \choice Clustering 2 has been obtained by single linkage.
  \choice Clustering 2 has been obtained by complete linkage.
  \choice Clustering 1 has been obtained by complete linkage.
  \choice One of the two clusterings could theoretically have been the result of either linkage function.
\end{itemize}
\closequestion

\questiontitle{Hierarchical clustering with average linkage on letters}
\repeatnote

You are given the set of data points
\[
x=\{a,c,e,h,i,m,n,s\},
\]
where the similarity between two points $s(x,x')$ is defined as the distance between the ordered letters in the alphabet,
for example $s(a,b)=1$, $s(a,c)=2$, and so on.
Perform agglomerative clustering using average linkage,
\[
d_{\mathrm{avg}}(A,B)=\frac{1}{n_A}\frac{1}{n_B}\sum_{a\in A}\sum_{b\in B}s(a,b),
\]
and stop when the distance between clusters is greater than $7$.
Answer the following questions with integer values.
\begin{itemize}
  \item How many clusters are formed in total? \answerline{2.0cm}
  \item What is the maximum number of data points in any one cluster? \answerline{2.0cm}
  \item What is the minimum distance between any two of the formed clusters? \answerline{2.0cm}
\end{itemize}
\closequestion

\questiontitle{Biclustering}
Which statements about Biclustering are correct?
Multiple answers are possible.
\begin{itemize}
  \choice A common biclustering method is the variance maximization method.
  \choice FABIA is a generative model.
  \choice FABIA defines biclusters as inner products.
  \choice In FABIA, the likelihood is analytically intractable and requires, for example, a variational EM algorithm.
  \choice Row and column elements must belong to exactly one bicluster or to no bicluster at all.
  \choice Biclustering simultaneously clusters the rows and columns of a matrix.
\end{itemize}
\closequestion

\newpage
\section{Dense Answer Key}

\begin{keybox}[Compact solutions]
\footnotesize
\begin{multicols}{2}
\begin{enumerate}[leftmargin=*,label=\textbf{Q\arabic*.}]
  \item $\kappa_4(X)=-0.13$, $\kappa_4(2X)=-2.13$, $\kappa_4(X+Y)=-0.27$, sub-Gaussian.
  \item $\kappa_3(X)=0$, $\kappa_4(X)=-2$, $\kappa_4(3X)/\kappa_4(X)=81$, $\kappa_4(X+Y)/\kappa_4(X)=2$.
  \item $A=12$.
  \item $D_{\mathrm{KL}}(p_Y\|p_Z)=5.00$.
  \item Third row $(12,15,18)$.
  \item $P(X=1)=0.63$, $P(Y=1)=0.50$, $H(X)=0.66$, $H(Y)=0.69$, $I(X;Y)=0.03$.
  \item $P(Y=0)=0.33$, $H(X)=0.64$, $H(Y)=0.64$, $I(X;Y)=0.17$.
  \item $P(X=0)=0.70$, $P(Y=0)=0.50$, $H(X)=0.61$, $H(Y)=0.69$, $I(X;Y)=0.02$.
  \item Correct: a, e, f, g.
  \item unbiased; variance; mean squared error; bias; biased but consistent; unbiased but not consistent; likelihood; density; variance; efficient.
  \item Correct: d, e, f, h.
  \item Correct: a, b, d, g, h.
  \item Correct: b, e, g, h.
  \item Correct: a, d, e, f, g.
  \item Correct: a, b, e, g.
  \item PCA corresponds to A, ICA corresponds to B; picture 1 is more likely PC2 vs. PC1 and picture 2 PC3 vs. PC1; in the EM part, $a$ is earlier and $b$ is later.
  \item $y_2=(3.46,0.00,0.00)$.
  \item $y_3=(4.24,0.00,0.00)$.
  \item Correct: c, d, e, f.
  \item Correct: b, d, h.
  \item Correct: a, b, c, d, f.
  \item Correct: c, e, h.
  \item Correct: a, b, d.
  \item Correct: a, d, g, k.
  \item First row $[10,11,17]$.
  \item Correct: a, b, c, e.
  \item $\mu^{\mathrm{new}}=(0,0)$, $\nu^{\mathrm{new}}=(3,1)$.
  \item Correct: a, d, e, h.
  \item $\mu^{\mathrm{new}}=(1,-1)$, $\nu^{\mathrm{new}}=(2.5,0.5)$.
  \item $\mu^{\mathrm{new}}=(0,-1)$, $\nu^{\mathrm{new}}=(3,1)$.
  \item Correct: b, d.
  \item Correct: a, b, c, f.
  \item Correct: c, d, e.
  \item Correct: a, b, f.
  \item Correct: a, c, d, e, f, g.
  \item Correct: a, d, e.
  \item $P(B\mid x=1.1)=0.25$.
  \item $P(j\mid x)=0.25$.
  \item $P(x\mid j)=0.60$.
  \item $P(G_1\mid x=1)\approx 0.66$.
  \item $D_{\mathrm{KL}}(p_Y\|p_Z)=6.00$.
  \item $D_{\mathrm{KL}}(p\|q)=0.19$.
  \item Correct: a, h.
  \item Correct: d, e, f, g, h.
  \item Correct: b, d, e.
  \item $3$ clusters; maximum cluster size $4$; minimum inter-cluster distance $8$.
  \item Correct: b, d, f.
\end{enumerate}
\end{multicols}
\end{keybox}

\end{document}
